% Assignment 1 - My first LaTeX document
% Tarun Konde (2025UCP1811)
%
% Changes made over the starter file:
%   - font size changed to 12pt and paper size set to a4 (default is letter)
%   - table of contents added
%   - an extra section and a subsection added
%   - numbered list
%   - "Science!" section with inline math
%   - natbib citation using physics.bib with the agu style
%   - labels and cross-references

\documentclass[12pt,a4paper]{article}

\usepackage[T1]{fontenc}
\usepackage{lmodern}

\usepackage{amsmath}
\usepackage{natbib}

\title{My First \LaTeX{} Document}
\author{Tarun Konde \\ 2025UCP1811}
\date{\today}

\begin{document}

\maketitle
\tableofcontents

\section{Introduction}\label{sec:intro}

This is my first document written in \LaTeX. Unlike a word processor, I write
plain text with commands and the compiler decides the layout.

Extra    spaces between words     are ignored. A single line break is also
ignored, so this line joins the previous one.

A blank line starts a new paragraph, like this one.

Some characters are reserved and have to be escaped to print them:
\# \$ \% \& \_ \{ \}. The backslash itself is printed with
\verb|\textbackslash|, which gives \textbackslash.

\subsection{Quotes and dashes}

Opening quotes are typed with backticks: ``double'' and `single'.
There are three kinds of dashes: a hyphen (well-known), an en dash
for ranges (pages 10--20) and an em dash for a break in a sentence---like this.

\section{Why I am learning \LaTeX}\label{sec:why}

\begin{enumerate}
    \item It handles equations far better than a word processor.
    \item Numbering of sections, figures and references is automatic.
    \item The source is plain text, so it works well with Git.
\end{enumerate}

Section~\ref{sec:intro} covered the basic rules of whitespace and special
characters.

\section{Science!}\label{sec:science}

In 1905 Einstein showed that mass and energy are equivalent, which is written
as $e = mc^2$ \citep{einstein1905}. Here $e$ is energy, $m$ is mass and $c$
is the speed of light in vacuum.

The same equation in display mode with a number:
\begin{equation}\label{eq:massenergy}
    e = mc^2
\end{equation}

Equation~\eqref{eq:massenergy} is discussed in Section~\ref{sec:science}
on page~\pageref{sec:science}. \citet{schrodinger1926} is another paper from
the same bib file, cited in text form.

\bibliographystyle{agu}
\bibliography{physics}

\end{document}
